A. IDENTITY

1. Essence is simple immediacy as sublated immediacy.
Its negativity is its being;
it is equal to itself in its absolute negativity
by virtue of which otherness and reference to other
have as such simply disappeared into pure self-equality.
Essence is therefore simple self-identity.

This self-identity is the immediacy of reflection.
It is not that self-equality which being is, or also nothing,
but a self-equality which, in producing itself as unity,
does not produce itself over again, as from another,
but is a pure production, from itself and in itself,
essential identity.

It is not, therefore, abstract identity
or an identity which is the result of
a relative negation preceding it,
one that separates indeed
what it distinguishes from it
but, for the rest, leaves it existing outside it,
the same after as before.

Being, and every determinateness of being,
has rather sublated itself not relatively,
but in itself, and this simple negativity,
the negativity of being in itself,
is the identity itself.

In general, therefore, it is still the same as essence.

B. DIFFERENCE

1. Absolute difference

Difference is the negativity that reflection possesses in itself,
the nothing which is said in identity discourse,
the essential moment of identity itself which,
as the negativity of itself, at the same time determines itself
and is differentiated from difference.

1. This difference is difference in and for itself,
absolute difference, the difference of essence.
It is difference in and for itself,
not difference through something external
but self-referring, hence simple, difference.
It is essential that we grasp absolute difference as simple.
In the absolute difference of A and not-A from each other,
it is the simple “not” which, as such, constitutes the difference.
Difference itself is a simple concept.
“In this,” so it is said, “two things differ, in that etc.”
“In this,” that is, in one and the same respect,
relative to the same basis of determination.
It is the difference of reflection, not the otherness of existence.
One existence and another existence are posited as lying outside each other;
each of the two existences thus determined over against each other
has an immediate being for itself.
The other of essence, by contrast, is the other in and for itself,
not the other of some other which is to be found outside it;
it is simple determinateness in itself.
Also in the sphere of existence did otherness and determinateness
prove to be of this nature, simple determinateness, identical opposition;
but this identity showed itself only as
the transition of a determinateness into the other.
Here, in the sphere of reflection, difference comes in
as reflected, so posited as it is in itself.

2. Difference in itself is the difference
that refers itself to itself;
thus it is the negativity of itself,
the difference not from another
but of itself from itself;
it is not itself but its other.
What is different from difference, however, is identity.
Difference is, therefore, itself and identity.
The two together constitute difference;
difference is the whole and its moment.
One can also say that difference,
as simple difference, is no difference;
it is such only with reference to identity;
even better, that as difference it entails
itself and this reference equally.
Difference is the whole and its own moment,
just as identity equally is its whole and its moment.
This is to be regarded as
the essential nature of reflection
and as the determined primordial origin
of all activity and self-movement.
Both difference and identity make themselves
into moment or positedness
because, as reflection, they are negative self-reference.
Difference, thus as unity of itself and of identity,
is internally determined difference.
It is not the transition into another,
not reference to another outside it;
it has its other, identity, within,
and in like manner identity,
in being included in the determination of difference,
has not lost itself in it as its other
but retains itself therein is
the reflection-into-itself of difference, its moment.

3. Difference has both these moments, identity and difference;
thus the two are both a positedness, determinateness.
But in this positedness each refers to itself.
The one, identity, is itself immediately
the moment of immanent reflection;
but no less is the other, difference,
difference in itself, reflected difference.
Difference, inasmuch as it has two such moments
which are themselves reflections into themselves,
is diversity.

2. Diversity

1. Identity internally breaks apart into diversity because, as absolute dif-
ference in itself, it posits itself as the negative of itself and these, its two
moments (itself and the negative of itself ), are reflections into themselves,
are identical with themselves; or precisely because it itself immediately sub-
lates its negating and is in its determination reflected into itself. The different
subsists as diverse, indifferent to any other, because it is identical with itself,
because identity constitutes its base and element; or, the diverse remains
what it is even in its opposite, identity.

Diversity constitutes the otherness as such of reflection. The other of
existence has immediate being, where negativity resides, for its foundation.
But in reflection it is self-identity, the reflected immediacy, that constitutes
the subsistence of the negative and its indifference.

The moments of difference are identity and difference itself. These
moments are diverse when reflected into themselves, referring themselves to
themselves; thus, in the determination of identity, they are only self-referring;
identity is not referred to difference, nor is difference referred to identity;
hence, inasmuch as each of these moments is referred only to itself, the two
are not determined with respect to each other.  Now because in this way
the two are not differentiated within, the difference is external to them. The
diverse moments, therefore, conduct themselves with respect to each other,
not as identity and difference, but only as moments different in general,
indifferent to each other and to their determinateness.

2. In diversity, as the indifference of difference, reflection has in general
become external; difference is only a positedness or as sublated, but is itself the
whole reflection.  On closer consideration, both, identity and difference
are reflections, as we have just established; each is the unity of it and
its other, each is the whole. But the determinateness, to be only identity
or only difference, is thus a sublated something. They are not, therefore,
qualities, since their determinateness, because of the immanent reflection,
is at the same time only as negation. What we have is therefore this
duplicity, immanent reflection as such and determinateness as negation or
positedness. Positedness is the reflection that is external to itself; it is negation
as negation and consequently, indeed in itself self-referring negation and
immanent reflection, but only in itself, implicitly; its reference is to a
something external.

Reflection in itself and external reflection are thus the two determi-
nations in which the moments of difference, identity and difference, are
posited. They are these moments themselves as they have determined them-
selves at this point.  Immanent reflection is identity, but determined to
be indifferent to difference, not to have difference at all but to conduct
itself towards difference as identical with itself; it is diversity. It is identity
that has so reflected itself into itself that it truly is the one reflection of
the two moments into themselves; both are immanent reflections. Identity
is this one reflection of the two, the identity which has difference within
it only as an indifferent difference and is diversity in general.  External
reflection, on the contrary, is their determinate difference, not as absolute
immanent reflection, but as a determination towards which the implicitly
present reflection is indifferent; its two moments, identity and difference
themselves, are thus externally posited, are not determinations that exist in
and for themselves.

Now this external identity is likeness,
and external difference is unlikeness.
Likeness is indeed identity, but only as a positedness,
an identity which is not in and for itself.
Unlikeness is equally difference,
but an external difference which is not, in and for itself,
the difference of the unlike itself.
Whether something is like or unlike something else is
not the concern of either the like or the unlike;
each refers only to itself, each is in and for itself what it is;
identity or non-identity, in the sense of likeness or unlikeness,
depend on the point of view of a third external to them.

3. External reflection connects diversity by referring it to likeness and
unlikeness. This reference, which is a comparing, moves back and forth from
likeness to unlikeness and from unlikeness to likeness. But this back and
forth referring of likeness and unlikeness is external to these determinations
themselves; moreover, they are not referred to each other, but each, for itself,
is referred to a third. In this alternation, each immediately stands out on its
own.  External reflection is as such external to itself; determinate difference
is negated absolute difference; it is not simple difference, therefore, not an
immanent reflection, but has this reflection outside it; hence its moments
come apart and both refer, each also outside the other, to the immanent
reflection confronting them.

In reflection thus alienated from itself, likeness and unlikeness present
themselves, therefore, as themselves unconnected, and reflection keeps them
apart, for it refers them to one and the same something by means of “in so
far,” “from this side or that,” and “from this view or that.” Thus diverse
things that are one and the same, when likeness and unlikeness are said of
them, are from one side like each other, but from another side unlike, and
in so far as they are alike, to that extent they are not unlike. Likeness thus
refers only to itself, and unlikeness is equally only unlikeness.

Because of this separation from each other, they sublate themselves.
Precisely that which should save them from contradiction and dissolution,
namely that something is like another in one respect but unlike in another
precisely this keeping of likeness and unlikeness apart, is their destruction.
For both are determinations of difference; they are references to each other,
each intended to be what the other is not; the like is not the unlike, and the
unlike is not the like; both have this connecting reference essentially, and
have no meaning outside it; as determinations of difference, each is what it
is as different from its other. But because of their indifference to each other,
the likeness is referred to itself, and similarly is unlikeness a point of view of
its own and a reflection unto itself; each, therefore, is like itself; difference
has vanished, since they have no determinateness to oppose them; in other
words, each is consequently only likeness.

Accordingly, this indifferent viewpoint or the external difference sublates
itself and it is in itself the negativity of itself. It is the negativity which in
comparing belongs to that which does the comparing. This latter oscillates
from likeness to unlikeness and back again; hence it lets the one disappear
into the other and is in fact the negative unity of both. This negative
unity transcends at first what is compared as well as the moments of the
comparing as a subjective operation that falls outside them. But the result
is that this unity is in fact the nature of likeness and unlikeness themselves.
Even the independent viewpoint that each of these is, is rather the self-
reference that sublates their distinctness and so, too, themselves.

From this side, as moments of external reflection and as external to
themselves, likeness and unlikeness disappear together into their likeness.
But this, their negative unity, is in addition also posited in them; for their
reflection implicitly exists outside them, that is, they are the likeness and
unlikeness of a third, of another than they themselves are. Thus the like is
not the like of itself, and the unlike, as the unlike not of itself but of an
unlike to it, is itself the like. The like and the unlike is each therefore the
unlike of itself. Each is thereby this reflection: likeness, that it is itself and
the unlikeness; unlikeness, that it is itself and the likeness.
Likeness and unlikeness constituted the side of positedness as against
what is being compared or the diverse which, as contrasted with them,
had determined itself as implicitly existent reflection. But this positedness
has consequently equally lost its determinateness as against this reflection.

Likeness and unlikeness, the determinations of external reflection, are pre-
cisely the merely implicitly existent reflection which the diverse as such was
supposed to be, its only indeterminate difference. Implicitly existent reflec-
tion is self-reference without negation, abstract self-identity and therefore
positedness itself.  The merely diverse thus passes over through the posit-
edness into negative reflection. The diverse is difference which is merely
posited, hence a difference which is no difference, hence a negation that
negates itself within. Likeness and unlikeness themselves, the positedness,
thus return through indifference or through implicitly existing reflection
back into negative unity with themselves, into the reflection which is the
implicit difference of likeness and unlikeness. Diversity, the indifferent sides
of which are just as much simply and solely moments of a negative unity, is
opposition.

3. Opposition

In opposition, the determinate reflection, difference, is brought to comple-
tion. Opposition is the unity of identity and diversity; its moments are
diverse in one identity, and so they are opposites.
Identity and difference are the moments of difference as held inside dif-
ference itself; they are reflected moments of its unity. Likeness and unlikeness
are instead the externalized reflection; their self-identity is not only the
indifference of each towards the other differentiated from it, but towards
being-in-and-for-itself as such; theirs is a self-identity that contrasts with
identity reflected into itself, hence an immediacy which is not reflected into
itself. The positedness of the sides of external reflection is therefore a being,
just as their non-positedness is a non-being.
On closer consideration, the moments of opposition are positedness
reflected into itself or determination in general. Positedness is likeness and
unlikeness; these two, reflected into themselves, constitute the determi-
nations of opposition. Their immanent reflection consists in that each is
within it the unity of likeness and unlikeness. Likeness is only in a reflec-
tion which compares according to the unlikeness and is therefore mediated
by its other indifferent moment; similarly, unlikeness is only in the same
reflective reference in which likeness is.  Each of these moments, in its
determinateness, is therefore the whole. It is the whole because it also con-
tains its other moment; but this, its other, is an indifferent existent; thus
each contains a reference to its non-being, and it is reflection-into-itself, or
the whole, only as essentially referring to its non-being.
This self-likeness, reflected into itself and containing the reference to
unlikeness within it, is the positive; and the unlikeness that contains within
itself the reference to its non-being, to likeness, is the negative.  Or again,
both are positedness; now in so far as the differentiated determinateness is
taken as a differentiated determinate reference of positedness to itself, oppo-
sition is, on the one hand, positedness reflected into its likeness with itself;
and, on the other hand, it is the same positedness reflected into its inequal-
ity with itself: the positive and the negative.  The positive is positedness as
reflected into self-likeness; but what is reflected is positedness, that is, the
negation as negation, and so this immanent reflection has the reference to
the other for its determination. The negative is positedness as reflected into
unlikeness; but positedness is the unlikeness itself, and so this reflection is
therefore the identity of unlikeness with itself and absolute self-reference.
Each, therefore, equally has the other in it: positedness reflected into self-
likeness has the unlikeness; and positedness reflected into self-unlikeness,
the likeness.

The positive and the negative are thus the sides of opposition that
have become self-subsisting. They are self-subsisting because they are the
reflection of the whole into itself, and they belong to opposition in so far
as the latter is determinateness which, as the whole, is reflected into itself.
Because of their self-subsistence, the opposition which they constitute is
implicitly determinate. Each is itself and its other; for this reason, each
has its determinateness not in an other but within.  Each refers itself to
itself only as referring itself to its other. This has a twofold aspect. Each
is the reference to its non-being as the sublating of this otherness in itself;
its non-being is thus only a moment in it. But, on the other hand, here
positedness has become a being, an indifferent subsistence; the other of
itself which each contains is therefore also the non-being of that in which
it should be contained only as a moment. Each is, therefore, only to the
extent that its non-being is, the two in an identical reference.
The determinations which constitute the positive and the negative con-
sist, therefore, in that the positive and the negative are, first, absolute
moments of opposition; their subsistence is indivisibly one reflection; it is
one mediation in which each is by virtue of the non-being of its other,
hence by virtue of its other or its own non-being.  Thus they are simply
opposites; or each is only the opposite of the other; the one is not yet the
positive and the other not yet the negative, but both are negative with
respect to each other. Each, therefore, simply is, first, to the extent that the
other is; it is what it is by virtue of the other, by virtue of its own non-being;
it is only positedness. Second, it is to the extent that the other is not; it is what
it is by virtue of the non-being of the other; it is reflection into itself.  The
two, however, are both the one mediation of opposition as such in which
they simply are only posited moments.

Moreover, this mere positedness is reflected into itself in general and,
according to this moment of external reflection, the positive and the negative
are indifferent towards this first identity where they are only moments; or
again, because that first reflection is the positive’s and the negative’s own
reflection into itself, each is indifferent towards its reflection into its non-
being, towards its own positedness. The two sides are thus merely diverse,
and because their determinateness  that they are positive or negative
constitutes their positedness as against each other, each is not specifically so
determined internally but is only determinateness in general; to each side,
therefore, there belongs indeed one of the two determinacies, the positive
or the negative; but the two can be interchanged, and each side is such as
can be taken equally as positive or negative.

But, in third place, the positive and the negative are not only a posited
being, nor are they something merely indifferent, but their positedness, or
the reference to the other in the one unity which they themselves are not, is
rather taken back into each. Each is itself positive and negative within; the
positive and the negative are the determination of reflection in and for
itself; only in this reflection of the opposite into itself is the opposite either
positive or negative. The positive has within it the reference to the other in
which the determinateness of the positive consists. And the same applies
to the negative: it is not negative as contrasted with another but has the
determinateness by which it is negative within.

Each is thus self-subsistent unity existing for itself. The positive is indeed
a positedness, but in such a way that the positedness is for it posited being
as sublated. It is the non-opposed, the sublated opposition, but as the side
of the opposition itself.  As positive, it is indeed a something which is
determined with reference to an otherness, but in such a way that its nature
is not to be something posited; it is the immanent reflection that negates
otherness. But its other, the negative, is itself no longer positedness or a
moment but itself a self-subsisting being; and so the negating reflection
of the positive is internally determined to exclude this being, which is its
non-being, from itself.

Thus the negative, as absolute reflection, is not the immediate negative
but is the negative as sublated positedness, the negative in and for itself
which positively rests upon itself. As immanent reflection, it negates its
reference to its other; its other is the positive, a self-subsisting being
hence its negative reference to this positive is the excluding of it from
itself. The negative is the independently existing opposite, over against the
positive which is the determination of the sublated opposition  the whole
opposition resting upon itself, opposed to the self-identical positedness.
The positive and the negative are such, therefore, not just in themselves,
but in and for themselves. They are in themselves positive and negative when
they are abstracted from their excluding reference to the other and are taken
only in accordance with their determination. Something is in itself positive
or negative when it is not supposed to be determined as positive or negative
merely in contrast with the other. But the positive and the negative, taken
not as a positedness and hence not as opposed, are each an immediate, being
and non-being. They are, however, moments of opposition: their in-itself
constitutes only the form of their immanent reflectedness. Something is
said to be positive in itself, outside the reference to something negative, and
something negative in itself, outside the reference to something negative: in
this determination, merely the abstract moment of this reflectedness is held
on to. However, to say that the positive and the negative exist in themselves
essentially implies that to be opposed is not a mere moment, nor that it is
just a matter of comparison, but that it is the determination of the sides
themselves of the opposition. The sides, as positive or negative in themselves,
are not, therefore, outside the reference to the other; on the contrary, this
reference, precisely as exclusive, constitutes their determination or their in-
itselfness; in this, therefore, they are at the same time in and for themselves.

C. CONTRADICTION

1. Difference in general contains both its sides as moments;
in diversity, these sides fall apart as indifferent to each other;
and in opposition as such, they are the moments of difference,
each determined by the other and hence only moments.
But in opposition these moments are equally determined within,
indifferent to each other and mutually exclusive, self-subsisting
determinations of reflection.

One is the positive and the other the negative, but the former as a positive
which is such within, and the latter as a negative which is such within. Each
has indifferent self-subsistence for itself by virtue of having the reference to
its other moment within it; each moment is thus the whole self-contained
opposition.  As this whole, each moment is self-mediated through its other
and contains this other. But it is also self-mediated through the non-being of
its other and is, therefore, a unity existing for itself and excluding the other
from itself.

Since the self-subsisting determination of reflection excludes the other
in the same respect as it contains it and is self-subsisting for precisely
this reason, in its self-subsistence the determination excludes its own self-
subsistence from itself. For this self-subsistence consists in that it contains
the determination which is other than it in itself and does not refer to
anything external for just this reason; but no less immediately in that it is
itself and excludes from itself the determination that negates it. And so it
is contradiction.

Difference as such is already implicitly contradiction; for it is the unity
of beings which are, only in so far as they are not one  and it is the
separation of beings which are, only in so far as they are separated in the
same reference connecting them. The positive and the negative, however, are
the posited contradiction, for, as negative unities, they are precisely their
self-positing and therein each the sublating of itself and the positing of
its opposite.  They constitute determining reflection as exclusive; for the
excluding is one act of distinguishing and each of the distinguished beings,
as exclusive, is itself the whole act of excluding, and so each excludes itself
internally.

If we look at the two self-subsisting determinations of reflection on
their own, the positive is positedness as reflected into likeness with itself
positedness which is not reference to another, hence subsistence inasmuch
as the positedness is sublated and excluded. But with this the positive makes
itself into the reference of a non-being  into a positedness.  In this way the
positive is contradiction  in that, as the positing of self-identity by the
excluding of the negative, it makes itself into a negative, hence into the other
which it excludes from itself. This last, as excluded, is posited free of the
one that excludes; hence, as reflected into itself and itself as excluding. The
reflection that excludes is thus the positing of the positive as excluding
the other, so that this positing immediately is the positing of its other
which excludes it.

This is the absolute contradiction of the positive; but it is immediately
the absolute contradiction of the negative; the positing of both in one
reflection.  Considered in itself as against the positive, the negative is
positedness as reflected into unlikeness to itself, the negative as negative. But
the negative is itself the unlike, the non-being of another; consequently,
reflection is in its unlikeness its reference rather to itself.  Negation in
general is the negative as quality or immediate determinateness; but taken
as negative, it is referred to the negative of itself, to its other. If this second
negative is taken only as identical with the first, then it is also only imme-
diate, just like the first; they are not taken, therefore, as each the other of
the other, hence not as negatives: the negative is not at all an immediate.
But now, since each is moreover equally the same as what the other is,
this reference connecting them as unequal is just as much their identical
connection.

This is therefore the same contradiction which the positive is, namely
positedness or negation as self-reference. But the positive is only implicitly
this contradiction, is contradiction only in itself; the negative, on the
contrary, is the posited contradiction; for in its reflection into itself, as a
negative which is in and for itself or a negative which is identical with
itself, its determination is to be the not-identical, the exclusion of identity.
The negative is this, to be identical with itself over against identity, and
consequently, because of this excluding reflection, to exclude itself from
itself.

The negative is therefore the whole opposition  the opposition which,
as opposition, rests upon itself; distinction that absolutely does not refer
itself to another; distinction which, as opposition, excludes identity from
itself, but thereby also excludes itself, for as reference to itself it determines
itself as the very identity which it excludes.

2. Contradiction resolves itself.

In the self-excluding reflection we have just considered, the positive and
the negative, each in its self-subsistence, sublates itself; each is simply the
passing over, or rather the self-translating of itself into its opposite. This
internal ceaseless vanishing of the opposites is the first unity that arises by
virtue of contradiction; it is the null.

But contradiction does not contain merely the negative; it also contains
the positive; or the self-excluding reflection is at the same time positing
reflection; the result of contradiction is not only the null.  The positive
and the negative constitute the positedness of the self-subsistence; their own
self-negation sublates it. It is this positedness which in truth founders to the
ground in contradiction.

The immanent reflection by virtue of which the sides of opposition are
turned into self-subsistent self-references is, first of all, their self-subsistence
as distinct moments; thus they are this self-subsistence only in themselves,
for they are still opposites, and that they are in themselves self-subsistent
constitutes their positedness. But their excluding reflection sublates this
positedness, turns them into self-subsistent beings existing in and for them-
selves, such as are self-subsistent not only in themselves but by virtue of their
negative reference to their other; in this way, their self-subsistence is also
posited. But, further, by thus being posited as self-subsistent, they make
themselves into a positedness. They fate themselves to founder, since they
determine themselves as self-identical, yet in their self-identity they are
rather the negative, a self-identity which is reference-to-other.

However, on closer examination, this excluding reflection is not only
this formal determination. It is self-subsistence existing in itself, and the
sublating of this positedness  is only through this sublating a unity that
exists for itself and is in fact self-subsistent. Of course, through the sublating
of otherness or positedness, positedness or the negative of an other is indeed
present again. But in fact, this negation is not just a return to the first
immediate reference to the other, is not positedness as sublated immediacy,
but positedness as sublated positedness. The excluding reflection of self-
subsistence, since it is excluding, makes itself a positedness but is just as
much the sublation of its positedness. It is sublating reference to itself; in
that reference, it first sublates the negative and it secondly posits itself as a
negative, and it is only this posited negative that it sublates; in sublating
the negative, it both posits and sublates it at the same time. In this way the
exclusive determination is itself that other of itself of which it is the negation;
the sublation of this positedness is not, therefore, once more positedness
as the negative of an other, but is self-withdrawal, positive self-unity. Self-
subsistence is thus unity that turns back into itself by virtue of its own
negation, for it turns into itself through the negation of its positedness. It
is the unity of essence  to be identical with itself through the negation not
of an other, but of itself.

3. According to this positive side, since self-subsistence in opposition,
as excluding reflection, makes itself into a positedness and equally sublates
this positedness, not only has opposition foundered but in foundering it has
gone back to its foundation, to its ground.  The excluding reflection of the
self-subsisting opposition turns it into a negative, something only posited;
it thereby reduces its formerly self-subsisting determinations, the positive
and the negative, to determinations which are only determinations; and the
positedness, since it is now made into positedness, has simply gone back
to its unity with itself; it is simple essence, but essence as ground. Through
the sublating of the determinations of essence, which are in themselves
self-contradictory, essence is restored, but restored in the determination of
an exclusive, reflective unity  a simple unity which determines itself as
negation, but in this positedness is immediately like itself and withdrawn
into itself.

In the first place, therefore, because of its contradiction, the self-
subsisting opposition goes back into a ground; this opposition is what comes
first, the immediate from which the beginning is made, while the sublated
opposition or the sublated positedness is itself a positedness. Accordingly,
essence is as ground a positedness, something that has become. But conversely,
only this has been posited, namely that the opposition or the posited-
ness is something sublated, only is as positedness. As ground, therefore,
essence is excluding reflection because it makes itself into a positedness;
because the opposition from which the start was just now made and
was the immediate is the merely posited determinate self-subsistence of
essence; because opposition only sublates itself within, whereas essence is
in its determinateness reflected into itself. As ground, therefore, essence
excludes itself from itself, it posits itself; its positedness  which is what
is excluded  is only as positedness, as identity of the negative with
itself. This self-subsistent is the negative posited as the negative, something
self-contradictory which, consequently, remains in the essence as in its
ground.

The resolved contradiction is therefore ground,
essence as unity of the positive and the negative.
In opposition, determinateness has progressed to self-subsistence;
but ground is this self-subsistence as completed;
in it, the negative is self-subsistent essence, but as negative;
and, as self-identical in this negativity,
ground is thus equally the positive.
In ground, therefore, opposition and its contradiction are
just as much removed as preserved.
Ground is essence as positive self-identity
which, however, at the same time
refers itself to itself as negativity
and therefore determines itself,
making itself into an excluded positedness;
but this positedness is the whole self-subsisting essence,
and essence is ground, self-identical in its negation and positive.
The self-contradictory self-subsistent opposition was itself,
therefore, already ground;
all that was added to it was the determination of self-unity
which emerges as each of the self-subsisting opposites
sublates itself and makes itself into its other,
thereby founders and sinks to the ground
but therein also reunites itself with itself;
thus in this foundering, that is, in its positedness or in the negation,
it rather is for the first time the essence
that is reflected into itself and self-identical.
